\documentclass[12pt]{article}
\usepackage[utf8]{inputenc}
\usepackage[T1]{fontenc}
\usepackage{lmodern}
\usepackage{microtype}
\usepackage[hidelinks]{hyperref}
\usepackage[a4paper, margin=3cm]{geometry}
\usepackage{parskip}

\title{Hello, World}
\author{Your Name}
\date{June 2026}

\begin{document}
\maketitle

This is the first post in the scriptorium. A place to write.

The word \textit{scriptorium} referred to the room in a medieval monastery where
monks copied manuscripts. Here it is something more modest: a folder on a hard
drive, a few text files, and the habit of returning to them.

\section*{What this is for}

Writing things down forces you to find out what you actually think. The act of
composing a sentence is not a record of a thought already formed — it is the
formation of the thought itself.\footnote{This is an old idea. Montaigne built
his whole project on it.}

That is reason enough to keep a place like this.

\end{document}
